\chapter{Literature Review}
\thispagestyle{empty}

\section{EEG and its frequency bands}
The EEG is the potential difference between scalp electrodes produced by the summed synchronous
activity of large groups of cortical neurons, with amplitudes of tens of microvolts. Electrodes
are placed according to the international 10--20 system: the letter gives the region (F frontal,
C central, T temporal, P parietal, O occipital), odd numbers lie over the left hemisphere, even
numbers over the right and ``z'' on the midline. Activity is conventionally described in the
delta (1--4~Hz), theta (4--8~Hz), alpha (8--13~Hz), beta (13--30~Hz) and gamma ($>$30~Hz) bands.
The recorded signal also contains non-cerebral contributions: eye movements produce large slow
deflections at frontal electrodes, muscle activity produces broadband activity in the beta and
gamma range, and mains interference appears at 50~Hz. A classifier uses any consistent
difference between groups, whatever its origin, so these contributions matter for detection.

\section{EEG findings in ADHD and IDD}
The most studied EEG finding in ADHD is increased theta and decreased beta power, summarised as
a raised theta/beta ratio~\cite{snyder2006}. A later meta-analysis found that the effect has
become smaller in more recent studies and cannot serve as a reliable diagnostic
measure~\cite{arns2013}. Stimulant medication shifts the EEG of children with ADHD towards that
of controls, mainly in the theta and beta bands~\cite{clarke2002}, so medication status must be
considered whenever ADHD and control groups are compared. EEG studies of intellectual disability
are fewer and the available recordings are small; the dataset used here~\cite{sareen2020data} was
published to support such work.

\section{Methods used for EEG classification}
\textbf{Adaptive decomposition.} EEG is non-stationary and contains several oscillatory
components at once. Empirical mode decomposition~\cite{huang1998} splits a signal into intrinsic
mode functions by an iterative algorithm; variational mode decomposition (VMD)~\cite{dragomiretskiy2014}
replaces the algorithm by an optimisation problem that finds $K$ band-limited modes; multivariate
VMD (MVMD)~\cite{rehman2019} extends it to several channels with centre frequencies shared by all
channels; and successive MVMD (SMVMD)~\cite{liu2022} extracts the modes one after another, so
that their number adapts to each signal.

\textbf{Hand-crafted features} include band powers from the power spectral
density~\cite{welch1967}, Hjorth parameters~\cite{hjorth1970}, entropies~\cite{richman2000,
bandt2002} and connectivity measures between channels such as the weighted phase-lag
index~\cite{vinck2011}. They are classified with $k$-nearest neighbours, support vector
machines~\cite{cortes1995}, random forests~\cite{breiman2001} or gradient-boosted
trees~\cite{chen2016, ke2017}.

\textbf{Riemannian methods} represent a window by the covariance matrix of its channels. These
matrices are symmetric positive definite and lie on a curved manifold; projecting them onto the
tangent space at their mean gives vectors that standard classifiers can
use~\cite{barachant2012, congedo2017}. Riemannian methods were developed for brain--computer
interfaces~\cite{congedo2017}; none of the studies on the two datasets reviewed below uses them.

\textbf{Deep learning} learns features from the raw signal, with compact convolutional networks
such as EEGNet~\cite{lawhern2018} and ShallowConvNet~\cite{schirrmeister2017}, recurrent networks
(LSTM) or hybrids of both~\cite{roy2019}.

\section{Published results on the two datasets}
Table~\ref{tab:lit} lists representative published results on the two datasets used in this
work, with the evaluation protocol \emph{as described by the authors}. Chandela
\emph{et al.}~\cite{chandela2025} detect both disorders with one method: SMVMD of each window,
nine statistical features per mode integrated with an energy weighting, mRMR feature
selection~\cite{peng2005} and $k$-nearest neighbours, with 10-fold cross-validation over windows.
Their accuracies of \R{PaperDSoneAcc}\% (IDD) and \R{PaperDStwoAcc}\% (ADHD) are the highest
reported for a single method covering both datasets, and this method is the starting point of the
present work. Several other ADHD studies on the same dataset report 99\% or more with channel
selection and $k$-NN~\cite{maheshwari2025}, a residual network with attention~\cite{bansal2025} or
CatBoost on a selected feature subset~\cite{mao2025}.

\begin{table}[H]
\caption{Published results on the two datasets, with the evaluation protocol as described in
each paper. ``Window-level'' means that the folds or the train/test split are formed from
windows or samples, so that data of the same person can appear in training and test; a
``not reported'' entry means that the paper does not state that subjects were kept apart.}
\label{tab:lit}\centering\small
\begin{adjustbox}{max width=\textwidth}
\begin{tabular}{p{3.3cm}p{1.2cm}p{4.4cm}p{3.9cm}r}
\hline
Study & Data & Method & Evaluation (as described) & Accuracy (\%) \\
\hline
Chandela \emph{et al.} 2025~\cite{chandela2025} & IDD & SMVMD features, mRMR, KNN & 10-fold over windows & \R{PaperDSoneAcc} \\
Chandela \emph{et al.} 2025~\cite{chandela2025} & ADHD & SMVMD features, mRMR, KNN & 10-fold over windows & \R{PaperDStwoAcc} \\
Maheshwari \emph{et al.} 2025~\cite{maheshwari2025} & ADHD & entropy-based channel selection, KNN & 10-fold / 70:30 split; subject separation not reported & 99.29 \\
Bansal \emph{et al.} 2025~\cite{bansal2025} & ADHD & autoencoder + ResNet with attention & 80:20 split of samples & 99.42 \\
Mao \emph{et al.} 2025~\cite{mao2025} & ADHD & CatBoost, 206 selected features & random 10-fold; subject separation not reported & 99.58 \\
\hline
Geethapriya and Pattabiraman 2026~\cite{geethapriya2026} & ADHD & SVM on spectral features & 10-fold over windows & 90.0 \\
 & & SVM on spectral features & \textbf{subject-wise} 10-fold & 65.6 \\
 & & LSTM & subject-wise 10-fold & 61.6 \\
 & & CNN--BiLSTM & subject-wise 10-fold & 65.7 \\
 & & CNN (their best) & subject-wise 10-fold & 72.7 \\
\hline
\end{tabular}
\end{adjustbox}
\end{table}

\section{Limitations of published evaluations}
\label{sec:leakage}
\textbf{Window-level splitting.} A recording is cut into many short, overlapping windows, and
the windows of one person share everything that is specific to that person and that recording:
head shape, electrode contact, artefacts and individual rhythms. If the folds are formed from
windows, the test windows of a person have near-identical neighbours in the training set, and a
classifier can score highly by recognising the person rather than the disorder. Saeb
\emph{et al.}~\cite{saeb2017} showed for clinical machine learning in general that such
record-wise evaluation overestimates the accuracy on new people, and Brookshire
\emph{et al.}~\cite{brookshire2024} found that most published deep-learning studies of
translational EEG are affected. For the ADHD dataset used here, Geethapriya and
Pattabiraman~\cite{geethapriya2026} measured the effect directly: the same SVM on spectral
features reaches 90.0\% with window-level folds and 65.6\% with subject-wise folds
(Table~\ref{tab:lit}). This project found the same for the starting method (Chapter~4): its
accuracy on ADHD windows falls from \R{DStwoSegWinAcc}\% to \R{DStwoSubjWinAcc}\% when the
children in the test set are new.

\textbf{Small samples.} With 14 participants (IDD) or 121 children (ADHD), cross-validation
estimates are uncertain by several percentage points~\cite{varoquaux2018}, so small differences
between methods are not meaningful.

\textbf{Heavy models.} The best subject-wise ADHD result found in the literature, 72.7\%, was
obtained with a convolutional network trained end to end~\cite{geethapriya2026}, and several of
the other studies use deep or recurrent networks. Such models are usually trained on a GPU and
are hard to interpret, which matters for a clinical aid and for deployment on simple hardware.

\textbf{Confounds.} In the ADHD dataset the children with ADHD were taking methylphenidate and
the controls were not (Chapter~3); published studies rarely discuss this, although stimulants
change the EEG~\cite{clarke2002}.

\section{Research gap}
None of the studies on these datasets reviewed here combines (i) a strict subject-wise
evaluation, (ii) one pipeline for both disorders, and (iii) a lightweight, interpretable model, and no subject-wise result on the IDD dataset was found. This project addresses
these points: it evaluates every method on unseen people, and proposes a pipeline built from
an adaptive decomposition, the Riemannian geometry of channel covariances and gradient-boosted
trees, which runs on a laptop CPU.
